\begin{lemma}[Observed closed-walk holonomy separates normalized densities]
    \label{lem:peps_wilson_density_obstruction}
    \lean{TNLean.PEPS.regularPhysicalCutMatrix_normalized_reduced_ne_of_closedWalk_holonomy,
        TNLean.PEPS.IsGIsometric.torusBond_normalized_density_ne_of_closedWalk_holonomy}\leanok
    Let a finite graph carry regular group-isometric tensors for a finite
    group, and let $R$ have connected induced graph. Fix a closed walk in
    that induced graph. If two actual bond insertions have different
    conjugacy classes of holonomy along this walk, their physical cuts
    $M,N$ on $R$ are nonzero and their trace-normalized row Gram matrices
    differ. The same conclusion holds for the actual native inserted-matrix
    torus cuts when both periods are at least three.
    The complement may be disconnected and no simple connectedness is required.
\end{lemma}
\begin{proof}\leanok
    \uses{thm:peps_closed_walk_class_physical_measurement,
        thm:peps_physical_cut_column_action,
        thm:peps_regular_twisted_state_nonzero}
    The complete observed closed-walk measurement has a projection $Q$
    selecting the first conjugacy class. Its uniform action on every actual
    open boundary column gives $QM=M$ and $QN=0$ after complementary contraction.
    Regular injectivity makes both coefficient vectors, and therefore both
    cuts, nonzero. If the normalized row Grams were equal, multiplying
    by $Q$ would give
    $\operatorname{tr}(MM^\dagger)^{-1}MM^\dagger=0$, contradicting $M\ne0$.
    The native coefficient cuts agree with the ordered-edge cuts entry by entry.
\end{proof}

\begin{definition}[An observed ring and a native flux crossing it]
    \label{def:peps_annular_flux_observation}
    \lean{TNLean.PEPS.torusAnnularObservationWalk,
        TNLean.PEPS.torusAnnularObservationRegion,
        TNLean.PEPS.torusAnnularFluxPath}\leanok
    On the seven-by-seven torus let $R$ be the twelve physical sites on
    the perimeter of the square with coordinates from zero to three.
    Let $w$ traverse this perimeter counterclockwise from $(0,0)$.
    Let $p$ be the four eastward native dual steps from $(1,1)$ to $(5,1)$.
    The former endpoint is inside the unobserved hole and the latter is
    outside the square.
\end{definition}

\begin{lemma}[An endpoint-free annular observation detects native flux]
    \label{lem:peps_annular_flux_density_obstruction}
    \lean{TNLean.PEPS.torusAnnularFluxPath_endpoint_corners_outside,
        TNLean.PEPS.IsGIsometric.torusAnnularFlux_normalized_density_ne_vacuum}\leanok
    Every physical corner of both endpoint plaquettes of $p$ lies outside
    $R$. For every finite group, nonidentity insertion $g$, and native
    regular group-isometric torus tensor, the actual cut of this flux path
    and the vacuum cut on $R$ are nonzero and have different normalized
    observed reduced densities.
\end{lemma}
\begin{proof}\leanok
    \uses{def:peps_annular_flux_observation,lem:peps_wilson_density_obstruction}
    The support of the perimeter walk induces a connected observed graph.
    The string has zero horizontal crossing counts and crosses exactly the
    four downward bonds based at $(2,1),(3,1),(4,1),(5,1)$.
    Only the second of these bonds belongs to the observed perimeter,
    where the upward transport is $g^{-1}$. Thus the actual observed loop
    holonomy is $g^{-1}$, while the vacuum holonomy is one.
    Their conjugacy classes differ because $g\ne1$. The physical closed-walk
    obstruction gives unequal normalized densities and both nonzero cuts.
    Direct inspection of the eight endpoint corners gives the stated exclusion.
\end{proof}

